\cotefile{69}
\scanpage{19}\mspage{15}
% @ top, lines 1–5 | text | structurel ; donnée de Q
\unsure{structurel|structural} : $\Psi$. Donc la donnée de \inserted{$Q_u$,} $Q$ avec cette \unsure{particularisation|restriction} équivaut à celle de $Q_u,P$. Pour reconstituer le graphe $C$, il faut reconstituer la notion de ``sommets liés'' sur $\{u\}\cup\underbrace{\left(\coprod_i P\wedge_{\Psi}(\mathbb F_2,p_i)\right)}_{E_{p}^*\ \deleted{\unsure{=Q_u}}}\amalg P$. On déclare :

% @ upper third, a | text | u est lié
a) $u$ est lié aux pts de $E_p^*$, et non à ceux de $P$

% @ upper third, b | text | un pt a de E
b) un \deleted{\unsure{est|un}} pt \inserted{$a$} de $E_p^*$ est lié à $\sigma a$ [i.e. si $a\in P\wedge_{\Psi}(\mathbb F_2,p_i)$, il est lié à l'autre pt $a'$ de cet ens] et à aucun autre pt de $E_p^*$

% @ upper middle, c | text | un pt a de E
c) un pt $a$ de $E_p^*$ est lié à un pt \deleted{\unsure{$a$}} $x$ de $P$ \unsure{ssi}, posant $a\in P\wedge_{\Psi}(\mathbb F_2,p_i)$, $a$ est dans l'image de $x$ par l'application canonique $P\longrightarrow P\wedge_{\Psi}(\mathbb F_2,p_i)$.

% @ middle, d | text | Deux pts x, y
d) Deux pts $x,y$ de $P$ sont liés sss l'élément $y-x=\alpha\in\Psi(Q_u)=\operatorname{Ker}(\mathbb F_{2}^{Q_u}\longrightarrow\mathbb F_2)$ est \deleted{\unsure{la}} de la forme \deleted{\unsure{$e_i$}} $e'_i$ $(i\in Q_u)$, où $e'_i=\sum_{j\in Q_u-\{i\}}e_j\in\Psi$ ($(e_j)_{j\in Q_u}$ étant la base canonique de $\mathbb F_2^{Q_u}$). Ceci ne fait que traduire le

% @ lower middle, lemma | text | Lemme Soient u ∈ C
\uline{Lemme} Soient $u\in C$, $x,y\in C-E(u)-\{u\}$. Pour que $x$ et $y$ soient liés, il \unsure{faut|f.} et s. qu'il \deleted{\unsure{l'ens|s'ensuive|puisse}} \deleted{\unsure{partie de l'ens}} \deleted{des 5 triangles $t$} \inserted{existe} \deleted{\uline{un triangle}} \inserted{(un et)} \uline{un seul} triangle de sommet $u$ tel que $t-E(x)=t-E(y)$ (NB les deux mbs sont \deleted{réduits} \inserted{de tte façon} à \uline{un} él. $\ne u$) soit \deleted{\unsure{réduis|réduits}} de cardinal 1) \unsure{SPS|SPS.|OPS} $x\ne y$.

% @ bottom, proof | text | Dém S'il n'existait
\uline{Dém} S'il n'existait aucun tel triangle, on aurait $x-y=\sum_{i\in\unsure{P_u|Q_u}}e_i$, or $\sum e_i\notin\Psi$, \inserted{car $5\ne0$ dans $\mathbb F_2$} absurde. Soit $\{u,v,w\}$ un tel triangle, \deleted{donc} $x,y$ sont liés tous deux à $v$ ou $w$ -- disons $v$. Si $y$ et $x$ sont liés, donc (avec les notations habituelles) $y=\sigma x$, on sait que pour tt $a\in E(u)^*\inserted{=E(u)-\{v,w\}}$, si $x$ lié à $a$, $\sigma x$ non lié à $a$ et inversement.
% archive sheet number: 18
